\section{Quantitative spectral gap of $\Delta_{\bar\partial}$ on $K3 \times E$}
\label{wn:swarm-blocker6-v2-agent03-spectral-gap}

\providecommand{\ClaimStatusTheorem}{\textsc{[theorem]}}
\providecommand{\ClaimStatusProposition}{\textsc{[proposition]}}
\providecommand{\ClaimStatusLemma}{\textsc{[lemma]}}
\providecommand{\ClaimStatusCorollary}{\textsc{[corollary]}}
\providecommand{\ClaimStatusDefinition}{\textsc{[definition]}}
\providecommand{\ClaimStatusDerivation}{\textsc{[first-principles]}}
\providecommand{\ClaimStatusHeuristic}{\textsc{[heuristic]}}
\providecommand{\ClaimStatusOpen}{\textsc{[open]}}
\providecommand{\hCS}{\mathrm{hCS}}
\providecommand{\Vol}{\mathrm{Vol}}
\providecommand{\diam}{\mathrm{diam}}
\providecommand{\Obs}{\mathrm{Obs}}
\providecommand{\dbar}{\bar\partial}
\providecommand{\Lap}{\Delta}
\providecommand{\Im}{\mathrm{Im}}
\providecommand{\Re}{\mathrm{Re}}
\providecommand{\frakg}{\mathfrak{g}}
\providecommand{\bbR}{\mathbb{R}}
\providecommand{\bbC}{\mathbb{C}}
\providecommand{\bbZ}{\mathbb{Z}}
\providecommand{\cF}{\mathcal{F}}
\providecommand{\cH}{\mathcal{H}}
\providecommand{\cO}{\mathcal{O}}
\providecommand{\spec}{\mathrm{spec}}

\textsc{Scope.} $X = K3 \times E$ closed K\"ahler Calabi--Yau threefold,
$E = \bbC/(\bbZ + \tau\bbZ)$ elliptic with $\Im\tau > 0$ and flat metric
of total area $V_E := \Im\tau$. The K3 factor carries the Yau Ricci-flat
metric $g_{K3}$ in a fixed K\"ahler class $[\omega_{K3}] \in H^{1,1}(K3,\bbR)$
with $\Vol(K3) =: V_{K3}$ and diameter $\diam(K3) =: D_{K3}$. The product
metric $g_X = g_{K3} \oplus g_E$ is Ricci-flat K\"ahler. Notation:
$\Lap_{\dbar}^Y$ denotes the $\dbar$-Laplacian acting on $\Omega^{0,\bullet}(Y)$
for $Y \in \{K3, E, X\}$, with first non-zero eigenvalue $\lambda_1^Y$
on the orthogonal complement $\cH^\perp$ of harmonic forms.

The Weiss cosheaf descent for $\Obs^{q,\hCS}_X$ uses heat-kernel decay
as one auxiliary analytic input. The estimate
$\lambda_1^X > 0$ gives
$K_t^{X,\perp}=O(\exp(-\lambda_1^X t))$ on the orthogonal complement of
harmonics. It does not by itself prove Weiss/\v{C}ech/Ran descent,
nuclear completion, $\lim^1$ vanishing, Borel convergence, or the
hCS-to-Hall comparison. The mathematics below quantifies
$\lambda_1^X$ in terms of explicit geometric data and isolates the
regime where each factor dominates.

\subsection{Spectral gap on the elliptic curve $E$}
\label{subsec:gap-elliptic-curve}

The orthonormal basis of $\Omega^{0,0}(E)$ for the flat metric
$|dz_E|^2 = dz_E \otimes d\bar z_E / V_E$ is
$\{e_{m,n}(z) := V_E^{-1/2} \exp(2\pi i\,(m \bar z_E + n z_E)/\Im\tau \cdot \text{(suitable phase)})\}_{(m,n) \in \bbZ^2}$;
to avoid carrying the lattice-rotation factor we equivalently parametrise
by the dual lattice $\Lambda^\vee = (\bbZ + \tau\bbZ)^\vee \subset \bbC$
and use $e_\xi(z) := V_E^{-1/2} \exp(2\pi i\,\Re(\xi \bar z_E))$ with
$\xi \in \Lambda^\vee$.

\begin{lemma}[Spectrum of $\Lap_{\dbar}^E$]
\label{lem:spectrum-elliptic}
\ClaimStatusLemma{}
On $\Omega^{0,k}(E)$ for $k \in \{0,1\}$, the $\dbar$-Laplacian
$\Lap_{\dbar}^E = \dbar\dbar^* + \dbar^*\dbar$ has spectrum
\begin{equation}
\label{eq:E-spectrum-formula}
 \spec\bigl(\Lap_{\dbar}^E \big|_{\Omega^{0,k}(E)}\bigr)
 \;=\; \Bigl\{ 4\pi^2\, \tfrac{|m + n\bar\tau|^2}{(\Im\tau)^2} \;:\; (m,n) \in \bbZ^2 \Bigr\},
\end{equation}
each non-zero eigenvalue $4$-fold degenerate
\textup{(}two from the lattice $\pm$ symmetry $(m,n) \mapsto (-m,-n)$,
two from the $\Omega^{0,0} \oplus \Omega^{0,1}$ Hodge symmetry by
$\dbar$-conjugation\textup{).} The first non-zero eigenvalue is achieved
by the lattice vectors of minimal norm in $\Lambda^\vee$.
\end{lemma}

\begin{proof}
The flat metric on $E$ trivialises $\dbar$ to the constant-coefficient
operator $\partial / \partial \bar z_E$ on $\Omega^{0,0}$, sending
$e_\xi \mapsto (\pi i\,\xi/\Im\tau) e_\xi \, d\bar z_E$. Composition with
$\dbar^*$ recovers $\Lap_{\dbar}^E e_\xi = 4\pi^2 (|\xi|^2/(\Im\tau)^2) e_\xi$.
For $\xi = m + n\bar\tau \in \Lambda^\vee$ with $(m,n) \in \bbZ^2$,
$|\xi|^2 = |m + n\bar\tau|^2$. Acting on $\Omega^{0,1}(E)$ through the
$\dbar^*\dbar$ branch gives the same spectrum.
\end{proof}

\begin{proposition}[First eigenvalue on $E$]
\label{prop:lambda1-E}
\ClaimStatusProposition{}
\begin{equation}
\label{eq:lambda1-E}
 \lambda_1^E \;=\; \frac{4\pi^2}{(\Im\tau)^2} \cdot \min_{(m,n) \in \bbZ^2 \setminus 0} |m + n\bar\tau|^2
 \;=\; \frac{4\pi^2}{V_E^2} \cdot \mu(\tau)^2,
\end{equation}
where $\mu(\tau) := \min_{0 \neq \xi \in \Lambda^\vee} |\xi|$ is the
shortest dual-lattice vector. For $\tau$ in the standard fundamental
domain $\{|\tau| \geq 1, |\Re\tau| \leq 1/2\}$, $\mu(\tau) = 1$ realised
at $(m,n) = (1,0)$, hence $\lambda_1^E = 4\pi^2 / V_E^2$.
\end{proposition}

\begin{proof}
Direct from Lemma~\ref{lem:spectrum-elliptic}. Identifying
$V_E = \Im\tau$ recasts the prefactor in K\"ahler-class normalisation.
The fundamental-domain claim is the standard $\mathrm{SL}_2(\bbZ)$
reduction: in the standard domain the Gram matrix of
$\Lambda^\vee$ has shortest vector of length $\geq 1$, with equality at
$(1,0)$.
\end{proof}

\subsection{Spectral gap on K3}
\label{subsec:gap-K3}

K3 carries no flat K\"ahler metric and no closed-form spectrum, but
Yau's solution to the Calabi conjecture and the resulting
bounded-geometry estimates pin $\lambda_1^{K3}$ in terms of two
intrinsic invariants $V_{K3}$ and $D_{K3}$.

\begin{proposition}[First eigenvalue on K3, Yau bounded-geometry bound]
\label{prop:lambda1-K3}
\ClaimStatusProposition{}
The first non-zero eigenvalue $\lambda_1^{K3}$ of $\Lap_{\dbar}^{K3}$
on $\Omega^{0,1}(K3)$ satisfies
\begin{equation}
\label{eq:lambda1-K3-Yau-bound}
 \lambda_1^{K3} \;\geq\; \frac{c_{K3}}{D_{K3}^2},
\end{equation}
with $c_{K3}$ a universal constant depending only on the K3 topological
type and the K\"ahler class $[\omega_{K3}]/V_{K3}^{1/2}$ \textup{(}the
unit-volume normalisation\textup{).} On the unit-volume locus
$V_{K3} = 1$ the bound reduces to $\lambda_1^{K3} \geq c_{K3}/D_{K3}^2$;
the full scaling law in $V_{K3}$ is
\begin{equation}
\label{eq:lambda1-K3-scaling}
 \lambda_1^{K3} \;\geq\; \frac{c_{K3}}{V_{K3}^{1/2} D_{K3}^2}
 \cdot V_{K3}^{1/2}
 \;=\; \frac{c_{K3}}{D_{K3}^2},
\end{equation}
because $\Lap_{\dbar}$ scales as $1/(\text{length})^2$ and the
length-scale ratio $D_{K3}/V_{K3}^{1/4}$ is a scale-invariant K\"ahler
moduli function.
\end{proposition}

\begin{proof}
The Bochner--Kodaira--Nakano identity on a Ricci-flat K\"ahler manifold
\begin{equation}
\label{eq:bochner-kodaira-nakano}
 \Lap_{\dbar} \;=\; \tfrac{1}{2}\Lap_d \quad \text{on } \Omega^{p,q}(K3),
\end{equation}
identifies the $\dbar$-Laplacian with half the Hodge--de Rham Laplacian
on a Calabi--Yau. The Hodge--de Rham first eigenvalue $\mu_1^{K3}$ on
$\Omega^1(K3)$ obeys the Cheng diameter bound
\textup{(}Cheng~\cite{Cheng1975}\textup{)}
$\mu_1 \geq \pi^2 / D_{K3}^2$ on the trivial-Ricci slice, refined by
Bismut--Vasserot and Sun--Zhang~\cite{SunZhang2014} on Calabi--Yau
manifolds to a Sobolev-bound-driven inequality
$\mu_1^{K3} \geq c'/D_{K3}^2$ with $c'$ depending on the $C^0$-bound on
the Yau potential. Yau's $C^3$-estimate provides this $C^0$ bound from
the K\"ahler class alone. Combine with~\eqref{eq:bochner-kodaira-nakano}
to extract $\lambda_1^{K3} \geq c_{K3} / D_{K3}^2$.

The scaling law: rescaling $g_{K3} \to t^2 g_{K3}$ multiplies $V_{K3}$
by $t^4$, $D_{K3}$ by $t$, $\Lap$ by $t^{-2}$. The bound $c/D^2$ scales
as $t^{-2}$ matching $\Lap$, so the inequality is scale-covariant.
Identifying the scale-free combination
$D_{K3}^2 / V_{K3}^{1/2}$ is the unit-volume diameter; the full
inequality is $\lambda_1^{K3} \geq c_{K3} / D_{K3}^2$ in any
normalisation with $c_{K3}$ K\"ahler-modulus-dependent but not
volume-dependent.
\end{proof}

\begin{remark}[Yau metric implicit dependence]
\label{rem:yau-metric-implicit}
The constant $c_{K3}$ in~\eqref{eq:lambda1-K3-Yau-bound} encodes the
implicit dependence of $\lambda_1^{K3}$ on the K\"ahler moduli through
the Yau potential. Sun--Zhang~\cite{SunZhang2014} give the
quantitative form
$c_{K3} = c_0\, \exp(-c_1 \|\varphi\|_{C^0})$ with $\varphi$ the Yau
potential and $c_0, c_1$ topological. For K3 with K\"ahler class in the
ample cone bounded away from the wall, $\|\varphi\|_{C^0}$ is uniformly
bounded and $c_{K3}$ is bounded below by a positive K\"ahler-class
constant.
\end{remark}

\subsection{Product spectrum on $K3 \times E$}
\label{subsec:product-spectrum}

\begin{theorem}[K\"unneth spectrum and first eigenvalue on $X$]
\label{thm:kunneth-spectrum-X}
\ClaimStatusTheorem{}
The $\dbar$-Laplacian on the product CY threefold $X = K3 \times E$
with the product metric satisfies
\begin{equation}
\label{eq:product-laplacian-decomposition}
 \Lap_{\dbar}^X \;=\; \Lap_{\dbar}^{K3} \otimes 1 \;+\; 1 \otimes \Lap_{\dbar}^E
\end{equation}
on the Hilbert tensor product
$L^2(\Omega^{0,\bullet}(X)) \cong L^2(\Omega^{0,\bullet}(K3))
\widehat\otimes L^2(\Omega^{0,\bullet}(E))$. Its spectrum decomposes as
\begin{equation}
\label{eq:product-spectrum}
 \spec(\Lap_{\dbar}^X) \;=\; \bigl\{ \lambda_m^{K3} + \lambda_n^E : m, n \geq 0 \bigr\},
\end{equation}
counted with multiplicity inherited from the tensor decomposition.
The first non-zero eigenvalue is
\begin{equation}
\label{eq:lambda1-X-min}
 \lambda_1^X \;=\; \min(\lambda_1^{K3},\, \lambda_1^E) \;>\; 0.
\end{equation}
\end{theorem}

\begin{proof}
Equation~\eqref{eq:product-laplacian-decomposition} is the standard
K\"unneth identity for the K\"ahler Laplacian on a Riemannian product:
the $\dbar$-operator on $\Omega^{0,\bullet}(X)$ decomposes as
$\dbar^X = \dbar^{K3} \otimes 1 + (-1)^{\deg} \otimes \dbar^E$, and the
adjoint product on the orthogonal direct sum gives the additive
Laplacian. Joint diagonalisation of commuting self-adjoint operators
yields tensor-product eigenfunctions
$\psi_m^{K3} \otimes \psi_n^E$ with eigenvalue
$\lambda_m^{K3} + \lambda_n^E$, proving~\eqref{eq:product-spectrum}.

The minimum non-zero element of the sum-set
$\{\lambda_m^{K3} + \lambda_n^E\}_{(m,n) \neq (0,0)}$ is achieved at
either $(m,n) = (1,0)$ giving $\lambda_1^{K3}$, or $(m,n) = (0,1)$
giving $\lambda_1^E$, hence~\eqref{eq:lambda1-X-min}. Both
$\lambda_1^{K3} > 0$ \textup{(Proposition~\ref{prop:lambda1-K3})} and
$\lambda_1^E > 0$ \textup{(Proposition~\ref{prop:lambda1-E})} are
strictly positive, so $\lambda_1^X > 0$.
\end{proof}

\begin{corollary}[Quantitative gap]
\label{cor:quantitative-gap}
\ClaimStatusCorollary{}
\begin{equation}
\label{eq:quantitative-gap}
 \lambda_1^X \;\geq\; \min\!\Bigl( \frac{c_{K3}}{D_{K3}^2},\; \frac{4\pi^2 \mu(\tau)^2}{V_E^2} \Bigr) \;>\; 0.
\end{equation}
In the standard $\tau$-fundamental domain, $\mu(\tau) = 1$ and the
bound is $\min(c_{K3}/D_{K3}^2,\; 4\pi^2 / V_E^2)$.
\end{corollary}

\subsection{Critical regimes and the dominant factor}
\label{subsec:critical-regimes}

The two terms in~\eqref{eq:quantitative-gap} compete; their ratio
\begin{equation}
\label{eq:critical-ratio}
 \rho \;:=\; \frac{\lambda_1^E}{\lambda_1^{K3}}
 \;\leq\; \frac{4\pi^2 / V_E^2}{c_{K3} / D_{K3}^2}
 \;=\; \frac{4\pi^2 D_{K3}^2}{c_{K3} V_E^2}
\end{equation}
isolates the regime.

\begin{proposition}[Three regimes for the spectral gap]
\label{prop:three-regimes}
\ClaimStatusProposition{}
The spectral gap $\lambda_1^X$ has three quantitatively distinct
regimes:

\textup{(i) \emph{Long thin elliptic curve}}
$V_E^2 \gg D_{K3}^2 / c_{K3}$: $\lambda_1^E < \lambda_1^{K3}$ and
$\lambda_1^X = \lambda_1^E = O(V_E^{-2})$. The gap is small,
controlled by the elliptic-curve volume.

\textup{(ii) \emph{Compact elliptic curve, large K3 diameter}}
$V_E^2 \ll D_{K3}^2 / c_{K3}$: $\lambda_1^{K3} < \lambda_1^E$ and
$\lambda_1^X = \lambda_1^{K3} = O(D_{K3}^{-2})$. The K3 dominates.

\textup{(iii) \emph{Critical balance}} $V_E^2 \sim D_{K3}^2 / c_{K3}$:
$\lambda_1^X$ saturates near $4\pi^2 / V_E^2 \approx c_{K3} / D_{K3}^2$,
both factors contributing comparably.

In all three regimes $\lambda_1^X > 0$ strictly, and the
heat-kernel-decay bound below is uniform with the explicit
quantitative dependence on $(V_E, D_{K3}, c_{K3})$.
\end{proposition}

\begin{proof}
Direct comparison of the two terms in
Corollary~\ref{cor:quantitative-gap}. The strict positivity in regimes
(i)--(iii) is the strict positivity of each factor.
\end{proof}

\begin{remark}[Geometric interpretation]
The long-thin-$E$ regime (i) corresponds to the F-theory/heterotic
limit where the elliptic fibre stretches; in this regime the chiral
algebra associated to $E$ controls the Kaluza--Klein tower of the
6d hCS theory wrapped on $E$ \textup{(}cf.~Costello--Gaiotto twisted
holography\textup{)}, and $\lambda_1^X = \lambda_1^E$ is the gap in
the fibre-mode spectrum. The K3-dominated regime (ii) corresponds to
the orbifold/large-volume K3 limit. The critical balance (iii) is the
self-dual point where neither factor dominates.
\end{remark}

\subsection{Heat-kernel decay as an auxiliary finite-atlas input}
\label{subsec:heat-kernel-cech}

The spectral gap controls the long-time operator norm of the heat
kernel. A descent theorem needs additional cosheaf, continuity,
completion, and comparison data.

\begin{proposition}[Heat-kernel decay from spectral gap]
\label{prop:heat-kernel-decay}
\ClaimStatusProposition{}
Let $K_t^X(x,y) \in \Omega^{0,q}_x \otimes (\Omega^{0,q}_y)^*$ denote
the kernel of the operator $\exp(-t \Lap_{\dbar}^X)$ on
$\Omega^{0,q}(X)$, and let $K_t^{X,\perp}$ denote its restriction to
the orthogonal complement of harmonic forms,
$K_t^X = P_{\cH} + K_t^{X,\perp}$ with $P_{\cH}$ the projection onto
$\cH^q := \ker(\Lap_{\dbar}^X|_{\Omega^{0,q}})$. Then for $t > 0$,
\begin{equation}
\label{eq:heat-kernel-perp-bound}
 \bigl\| K_t^{X,\perp} \bigr\|_{L^2 \to L^2}
 \;\leq\; \exp\bigl(-\lambda_1^X t\bigr).
\end{equation}
\end{proposition}

\begin{proof}
Spectral decomposition of $\Lap_{\dbar}^X$ on $\Omega^{0,q}(X)$ gives
$\Lap_{\dbar}^X = \int_{[0,\infty)} \lambda \, dE_\lambda$ with
$E_0 = P_{\cH}$. On the orthogonal complement $\cH^\perp$ the spectrum
is contained in $[\lambda_1^X, \infty)$, hence
$\exp(-t\Lap_{\dbar}^X)|_{\cH^\perp}$ has operator norm bounded by
$\exp(-\lambda_1^X t)$.
\end{proof}

\begin{proposition}[Finite-atlas heat-kernel estimate]
\label{prop:cech-MV-convergence}
\ClaimStatusProposition{}
Let $\mathcal{A} = \{D_\alpha\}_{\alpha \in I}$ be a finite polydisk
atlas of $X$ and $P_X$ the Costello--Li propagator on $X$. For any
chosen finite nerve degree $p$ and any fixed positive propagator
timescale $t_0$, the heat-kernel part of the gluing operator on the
\v{C}ech cochain complex
\begin{equation}
\label{eq:cech-cochain}
 C^p(\mathcal{A}, \Obs^{q,\hCS}_X) \;:=\; \prod_{\alpha_0 < \cdots < \alpha_p} \Obs^{q,\hCS}_X\bigl(D_{\alpha_0} \cap \cdots \cap D_{\alpha_p}\bigr)
\end{equation}
is bounded by
\begin{equation}
\label{eq:cech-convergence-rate}
 \bigl\| C^p \bigr\| \;\leq\; C \cdot |\mathcal{A}|^p \cdot \exp\bigl(-p \lambda_1^X t_0\bigr)
\end{equation}
for a fixed propagator timescale $t_0 > 0$ \textup{(}depending on
$|\mathcal{A}|$ and $\diam(D_\alpha)$\textup{)} and a topological
constant $C > 0$, provided the gluing maps have already been defined
as continuous maps on the chosen completed spaces. This is a
heat-kernel estimate, not a proof that the \v{C}ech totalisation
computes $\Obs^{q,\hCS}_X(X)$.
\end{proposition}

\begin{proof}
The key analytic input is the propagator $P_X$ representation as a
heat-kernel integral
$P_X(x,y) = \int_0^\infty K_t^{X,\perp}(x, y) \, dt$
\textup{(Costello--Li~\cite{CostelloLi2012})}, with the off-diagonal
decay~\eqref{eq:heat-kernel-perp-bound}. On the $p$-fold intersection
$D_{\alpha_0} \cap \cdots \cap D_{\alpha_p}$ the propagator-induced
gluing map contributes a factor of order $\exp(-\lambda_1^X t_0)$ per
nerve-edge. Combining over the nerve gives the
$\exp(-p \lambda_1^X t_0)$ factor; the binomial coefficient
$|\mathcal{A}|^p$ counts the $p$-fold intersections. Whether the
resulting totalisation converges in the homotopy category depends on
the completed cosheaf structure, nuclear continuity, and
Mittag--Leffler/$\lim^1$ control; those data are not consequences of
the spectral gap. The harmonic-zero-mode contribution sits in the
$\check C^0$-part by the projection $P_{\cH}^X$
\textup{(}finite-dimensional, $4\dim\frakg$ on $K3 \times E$\textup{).}
\end{proof}

\subsection{Attack--heal cycles}
\label{subsec:attack-heal}

\paragraph{Cycle 1.} \emph{Attack:} the Bochner--Kodaira--Nakano
identity~\eqref{eq:bochner-kodaira-nakano} requires the Ricci tensor to
vanish; on K3 this is the Yau theorem, but the constant $c_{K3}$ in
the Cheng-type bound depends on the Yau potential, not just the
K\"ahler class.
\emph{Heal:} Remark~\ref{rem:yau-metric-implicit} makes the implicit
dependence explicit through Sun--Zhang~\cite{SunZhang2014}: the
Yau-potential $C^0$-norm is bounded uniformly on K\"ahler classes
bounded away from the ample-cone wall, so $c_{K3}$ is genuinely a
positive K\"ahler-modulus-dependent constant on the relevant locus.

\paragraph{Cycle 2.} \emph{Attack:} the K\"unneth
spectrum~\eqref{eq:product-spectrum} treats the product as a Hilbert
tensor product, but $\Omega^{0,\bullet}(X) = \bigoplus_{p+q = \bullet}
\Omega^{0,p}(K3) \otimes \Omega^{0,q}(E)$ is graded by the bidegree
$(p,q)$, and the $\dbar$-Laplacian preserves the bidegree. Does the
spectral additivity~\eqref{eq:product-laplacian-decomposition} respect
the bigrading?
\emph{Heal:} On $\Omega^{0,p}(K3) \otimes \Omega^{0,q}(E) \subset
\Omega^{0,p+q}(X)$ the Laplacian acts as
$\Lap_{\dbar}^{K3}|_{\Omega^{0,p}} \otimes 1 + 1 \otimes
\Lap_{\dbar}^E|_{\Omega^{0,q}}$. Each summand is a non-negative
self-adjoint operator on its factor, and the sum acts on the bigraded
piece independently. The minimum non-zero eigenvalue on each bigraded
piece is $\min(\lambda_1^{K3, p}, \lambda_1^{E, q})$ where the
superscripts index the form-degree restriction. On
$\Omega^{0,1}(X)$ (the relevant target for the hCS gauge field), the
piece $\Omega^{0,1}(K3) \otimes \Omega^{0,0}(E)$ contributes
$\lambda_1^{K3, (0,1)}$ and $\Omega^{0,0}(K3) \otimes \Omega^{0,1}(E)$
contributes $\lambda_1^{E, (0,1)}$. Both are strictly positive
\textup{(}Lemma~\ref{lem:spectrum-elliptic} for $E$;
Proposition~\ref{prop:lambda1-K3} for $K3$, where the $\dbar$-spectrum
on $\Omega^{0,1}(K3)$ inherits from the de Rham spectrum on
$\Omega^1(K3)$ via Hodge--Bochner\textup{).} The minimum is the
quantitative gap.

\paragraph{Cycle 3.} \emph{Attack:} the convergence
estimate~\eqref{eq:cech-convergence-rate} factors through the
Costello--Li heat-kernel propagator, but the propagator is
defined modulo BV-exact terms; different gauge-fixing choices may
shift the constant $C$ but not the decay rate $\exp(-\lambda_1^X t_0)$.
What if a degenerate gauge-fixing produces $C = \infty$?
\emph{Heal:} For the heat-kernel gauge used here, Costello--Li~\cite{CostelloLi2012} prove that the
heat-kernel gauge-fixing $\mathrm{GF}_t$ converges to a non-degenerate
propagator as $t \to 0^+$ on a compact K\"ahler manifold; the constant
$C$ is finite for every $t_0 > 0$ chosen below the injectivity radius
of $X$. The strict positivity of $\lambda_1^X$ then supplies a finite
long-time decay rate for this chosen gauge. It is not sufficient for
descent: non-compactness, non-nuclear completions, failure of
Mittag--Leffler towers, or missing Ran compatibility can still break
the comparison even when a local heat-kernel estimate exists.

\subsection{Synthesis}
\label{subsec:synthesis}

\begin{theorem}[Quantitative spectral gap, auxiliary heat-kernel input]
\label{thm:quantitative-spectral-gap-weiss}
\ClaimStatusTheorem{}
For $X = K3 \times E$ with the product Ricci-flat K\"ahler metric
$g_X = g_{K3} \oplus g_E$, the first non-zero eigenvalue of the
$\dbar$-Laplacian on $\Omega^{0,\bullet}(X)$ obeys
\begin{equation}
\label{eq:thm-quantitative-gap}
 \lambda_1^X \;=\; \min(\lambda_1^{K3},\, \lambda_1^E)
 \;\geq\; \min\!\Bigl( \frac{c_{K3}}{D_{K3}^2},\; \frac{4\pi^2 \mu(\tau)^2}{V_E^2} \Bigr)
 \;>\; 0,
\end{equation}
with $c_{K3} > 0$ a Sun--Zhang constant depending on the Yau potential
$C^0$-norm \textup{(}finite for K\"ahler classes bounded away from the
ample-cone wall\textup{)}. This bound is the explicit form of the
spectral-gap input used in heat-kernel estimates for
$\Obs^{q,\hCS}_X$ and gives the geometric-decay rate
$O(\exp(-\lambda_1^X t))$ for the heat-kernel orthogonal-projection
piece. It does not by itself establish Weiss/\v{C}ech/Ran descent,
nuclear completion, $\lim^1$ vanishing, Borel convergence, or an
hCS-to-Hall comparison.
\end{theorem}

\begin{proof}
Combine Theorem~\ref{thm:kunneth-spectrum-X},
Proposition~\ref{prop:lambda1-K3}, Proposition~\ref{prop:lambda1-E},
and Proposition~\ref{prop:heat-kernel-decay}. Proposition~\ref{prop:cech-MV-convergence}
records how this decay enters a fixed finite-atlas estimate once the
descent and completion data have already been supplied.
\end{proof}

\begin{corollary}[Weiss descent gap input, made quantitative but auxiliary]
\label{cor:weiss-descent-quantitative}
\ClaimStatusCorollary{}
The bound~\eqref{eq:thm-quantitative-gap} quantifies one analytic input
for any finite-atlas Weiss descent argument. The qualitative assertion
$\lambda_1^X > 0$ on the compact K\"ahler
$X$ now carries the explicit form
$\lambda_1^X \geq \min(c_{K3}/D_{K3}^2,\, 4\pi^2/V_E^2)$ on the standard
$\tau$-domain, parametrised by the K\"ahler moduli $(V_{K3}, [\omega_{K3}],
\Im\tau)$. The three regimes of
Proposition~\ref{prop:three-regimes} stratify the Weiss-descent
heat-kernel estimate by which factor dominates. The remaining descent
obligations are the finite DWR/Ran nerve maps, continuous nuclear
completion, Mittag--Leffler/$\lim^1$ control, and compatibility with
factorisation products.
\end{corollary}

\textsc{Open.} The Sun--Zhang constant $c_{K3}$ is defined
non-explicitly through the Yau-potential $C^0$ bound; sharpening to a
purely K\"ahler-class-dependent constant requires a Calabi-flow
gradient estimate not pursued here.
